\graphicspath{{images/}{../images/},{../../images/}}

\setlrmarginsandblock{1in}{1in}{*}
\setulmarginsandblock{1in}{*}{1}
\checkandfixthelayout 

\hypersetup{
    %bookmarks=true,                             % show bookmarks bar?
    pdftitle={Estimating the sampling distribution of the maximum likelihood estimator of the parameters of a series system given a sample of masked system failure times},  % title
    pdfauthor={Alexander Towell},               % author
    pdfsubject={statistical inference},                    % subject of the document
    pdfkeywords={maximum likelihood estimation,
        information matrix,
        covariance matrix,
        multivariate normal,
        sampling distribution,
        statistical models,
        statistical inference,
        parameter estimation,
        parametric bootstrap,
        entropy,
        maximum entropy distribution,
        series systems,
        m-out-of-m system,
        exponential distribution,
        masked system failure time},           % keywords
    colorlinks=true,                            % false: boxed links; true: colored links
    linkcolor=blue,                             % color of internal links
    citecolor=green,                            % color of links to bibliography
    filecolor=magenta,                          % color of file links
    urlcolor=green                              % color of external links
}

\definecolor{lightgray}{RGB}{2, 2, 2}
\usepgfplotslibrary{fillbetween}
\pgfplotsset{compat=1.13}

\DeclareMathOperator{\likelihood}{\mathcal{L}}
\DeclareMathOperator{\expectation}{\mathbb{E}}
\DeclareMathOperator*{\argmax}{arg\,max}
\DeclareMathOperator*{\argmin}{arg\,min}
\DeclareMathOperator{\hessian}{\mathcal{H}}
\DeclareMathOperator{\jacobian}{\mathcal{J}}
\DeclareMathOperator{\powerset}{\mathcal{P}}
\DeclareMathOperator{\infomatrx}{\matrx{\mathcal{I}}}
\DeclareMathOperator{\infomatrxn}{\matrx{\mathcal{I}_n}}
\DeclareMathOperator{\observed}{\matrx{\mathcal{J}}}
\DeclareMathOperator{\observedn}{\matrx{\mathcal{J}_n}}
\DeclareMathOperator{\orderfunc}{\pi}
\DeclareMathOperator{\score}{\rm{s}}
\DeclareMathOperator{\entropy}{H}
\DeclareMathOperator{\info}{I}
\DeclareMathOperator{\mvn}{\rm{MVN}}
\DeclareMathOperator{\expdist}{\rm{EXP}}
\DeclareMathOperator{\var}{\rm{Var}}
\DeclareMathOperator{\infi}{\rm{inf}}
\DeclareMathOperator{\mse}{\rm{MSE}}
\DeclareMathOperator{\cov}{\matrx{\mathlarger{\Sigma_n}}}
\DeclareMathOperator{\trace}{tr}
\DeclareMathOperator{\bias}{bias}
\DeclareMathOperator{\scalarstat}{g}

\renewcommand{\vec}[1]{\boldsymbol{\mathbf{#1}}}

\newcommand{\covn}[1]{\operatorname{\matrx{\mathlarger{\Sigma_{#1}}}}}
\newcommand{\pfam}{\mathcal{F}_{\sysparam}}
\newcommand{\pfami}{\mathcal{F}_{\cparamv}}
\newcommand{\infomatrxni}[1]{\operatorname{\matrx{\mathcal{I}_{#1}}}}
\newcommand{\cardinality}[1]{{\lvert#1\rvert}}
\newcommand{\given}{{\,|\,}}
\newcommand{\matrx}[1]{\boldsymbol{#1}}
\newcommand{\vtopr}{\operatorname{vec}}
\newcommand{\vt}[1]{\vtopr\left(#1\right)}
\newcommand{\vectomat}[1]{\operatorname{mat}(#1)}
\newcommand{\biconditional}{\leftrightarrow}
\newcommand{\indicator}[1]{\operatorname{\mathbbm{1}_{#1}}}
\newcommand{\cntsymb}{\boldsymbol{\hat\omega}}
\newcommand{\cnt}[1]{\cntsymb_{\{#1\}}}
\newcommand{\veclist}[1]{ \left (#1 \right )}
\newcommand{\derivative}[1]{\frac {\mathrm{d}}{\mathrm{d}#1}}
\newcommand{\pderivative}[1]{\frac {\mathrm{\partial}}{\mathrm{\partial}#1}}
\newcommand{\ppderivative}[2]{\frac {\mathrm{\partial} #1}{\mathrm{\partial}#2}}
\newcommand{\psecondderivative}[2]{\frac {\mathrm{\partial^2}}{\mathrm{\partial}#1\mathrm{\partial}#2}}
\newcommand{\psecondderivativesame}[1]{\frac {\mathrm{\partial^2}}{\mathrm{\partial}#1^2}}
\newcommand{\rv}[1]{\mathrm{#1}}
\newcommand{\cdf}{F}
\newcommand{\pdf}{f}
\newcommand{\haz}{h}
\newcommand{\srv}{R}
\newcommand{\pmf}{p}
\newcommand{\prob}{\Pr} %\mathrm{Pr}
\newcommand{\kth}{$k^{\text{th }}$}
\newcommand{\ith}{$i^{\text{th }}$}
\newcommand{\jth}{$j^{\text{th }}$}
\newcommand{\nth}{$n^{\text{th }}$}
\newcommand{\rth}{$r^{\text{th }}$}
\newcommand{\pth}{$p^{\text{th }}$}
\newcommand{\estimator}[1]{\hat{#1}}
\newcommand{\nonmleestimator}[1]{\bar{#1}} %\tidle{#1}
\newcommand{\tparam}[1]{{#1}^{\star}}
\newcommand{\vecstat}{\vec{\scalarstat}}

\newcount\colveccount
\newcommand*\colvec[1]{
        \global\colveccount#1
        \begin{pmatrix}
        \colvecnext
}
\def\colvecnext#1{
        #1
        \global\advance\colveccount-1
        \ifnum\colveccount>0
                \\
                \expandafter\colvecnext
        \else
                \end{pmatrix}
        \fi
}

\makeatletter
\def\namedlabel#1#2{\begingroup
    #2%
    \def\@currentlabel{#2}%
    \phantomsection\label{#1}\endgroup
}
\makeatother

\newcommand{\cparam}{v}
\newcommand{\param}{\Theta}
\newcommand{\paramv}{\theta}
\newcommand{\sysparam}{\matrx{\param}}
\newcommand{\sysparamvec}{\vec{\paramv}}

\newcommand{\cparamv}{\vec{\cparam}}

\newcommand{\mfswn}[2]{\vec{M}(#2,#1)}
\newcommand{\mfsn}[1]{\vec{M_{#1}}}
\newcommand{\mf}{\vec{M}}
\newcommand{\mfo}{\vec{m}}
\newcommand{\mfs}{\vec{M_n}}
\newcommand{\mfsw}{\vec{M}(w,n)}

\newcommand{\mfsi}[1]{\vec{M}(w,{#1})}
\newcommand{\mfso}{\vec{m_n}}
\newcommand{\mfsoi}[1]{\vec{m_{#1}}}
\newcommand{\mfsoir}[2]{\vec{m_{#1}}^{#2}}

\newcommand{\pointest}{\vec{\nonmleestimator\sysparamvec_{n}}}
\newcommand{\pointesti}[1]{\vec{\nonmleestimator\sysparamvec_{n}^{(#1)}}}
\newcommand{\sampdpt}{\rv{\vec{Y_n}}}

\newcommand{\mleir}[2]{\vec{\estimator{\sysparamvec}}(#1 \given #2)}
\newcommand{\mlei}[1]{\vec{\estimator{\sysparamvec}}^{(#1)}(n \given w)}
\newcommand{\mleu}{\vec{\estimator{\sysparamvec}_{n}}}
\newcommand{\mle}{\vec{\estimator{\sysparamvec}}(n \given w)}
\newcommand{\mleui}[1]{\vec{\estimator{\sysparamvec}_{n}^{(#1)}}}

\newcommand{\expomle}{\vec{\estimator{\vec\lambda}_n}}
\newcommand{\expomlei}[2]{\vec{\estimator{\vec\lambda}_{#2}^{(#1)}}}

\newcommand{\sampdu}{\rv{\vec{Y_{\!n}}}}
\newcommand{\sampd}{\rv{\vec{Y}}(n \given w)}
\newcommand{\sampdi}[1]{\rv{\vec{Y_{\!n}^{(#1)}}}}
\newcommand{\sampdir}[3]{\rv{\vec{Y}}^{(#3)}(#1 \given #2)}
\newcommand{\sampdexpo}{\rv{\vec{\Lambda_n}}}
\newcommand{\sampdexpoi}[2]{\rv{\vec{\Lambda_{#2}}^{(#1)}}}

\newcommand{\cands}{\boldsymbol{\mathcal{C}}}
\newcommand{\candsw}[1]{\boldsymbol{\mathcal{C}_{#1}}}
\newcommand{\cand}{\vec{\rm{c}}}
\newcommand{\rvcand}{\rv{\vec{\rm{C}}}}
\newcommand{\rvcandi}[1]{\rv{\vec{\rm{C}_{#1}}}}
\newcommand{\cv}{\rv{T}}
\newcommand{\sv}{\rv{S}}
\newcommand{\ti}{\rm{t}}

\newcommand{\harmonic}[2]{H_{#1,#2}}

\newcommand{\spdf}{\pdf_{\sv}}
\newcommand{\jointpdf}{\pdf_{\rvcand, \sv \given \rv{W}}}

\numberwithin{equation}{chapter}

\theoremstyle{plain}
\newtheorem{theorem}{Theorem}[chapter]
\newtheorem{corollary}{Corollary}[theorem]

\newtheorem{definition}{Definition}[chapter]
\newtheorem{postulate}{Postulate}[chapter]
\newtheorem{assumption}{Assumption}[chapter]

\theoremstyle{definition}
\newtheorem{example}{Example}

\theoremstyle{remark}
\newtheorem*{remark}{Remark}
\newtheorem*{notation}{Notation}

\newcounter{equationstore}
\newcounter{proofnum}

\AtBeginEnvironment{proof}{\setcounter{equationstore}{\value{equation}}
\refstepcounter{proofnum}
\setcounter{equation}{0}\renewcommand{\theequation}{\alph{equation}}
\renewcommand{\theHequation}{\arabic{proofnum}.\alph{equation}}}
\AtEndEnvironment{proof}{\setcounter{equation}{\value{equationstore}}}

\SetKwFunction{findmle}{FindMLE}
\SetKwFunction{proj}{Project}
\SetKwFunction{genSampleFn}{GenerateMSFT}
\SetKwFunction{genSampleMLE}{GenerateMLE}